%
% 32-integration.tex -- Umlaufzahl und komplexe Integration
%
% (c) 2026 Prof Dr Andreas Müller
%
\subsection{Umlaufzahl und komplexe Integration}
Die Umlaufzahl 
nach Definition~\ref{buch:integration:definition:umlaufzahl}
zu berechnen ist eher mühsam.
Daher ist es besonders erfreulich, dass man sie auch mit komplexer
Wegintegration berechnen kann, wie der folgende Satz zeigt.
Der Beweis beruht auf der Idee, dass sich der in der
Abbildung~\ref{buch:integration:fig:argument} rot gezeichnete Weg
in den grünen Weg deformieren lässt, der die Richtung $s$ immer in der
gleichen Richtung überschreitet und damit seinerseits homotop zu
einem Weg der Form $t\mapsto e^{2\pi int}$, für den sich das Integral
leicht berechnen lässt.

\begin{satz}
Sei $\gamma$ ein Weg in $\mathbb{C}$, der den Punkt $z_0$ nicht trifft,
dann umläuft der Weg den Punkt $z_0$
\[
\operatorname{ind}_\gamma(z_0)
=
\frac{1}{2\pi i}
\oint_\gamma \frac{dz}{z-z_0}
\]
mal.
\end{satz}

\begin{proof}
Abbildung~\ref{buch:integration:fig:argument} zeigt, wie die Umlaufzahl
berechnet wird.
Durch ``Strecken'' des Graphen von $\varphi(t)$ zu einer Geraden
(in der Abbildung grün eingezeichnet), kann der Weg durch ein
Homotopie in einen Weg deformiert werden, der durch  die
Parametrisierung $\gamma_1(t) = z_0+re^{2\pi i n t}$ beschrieben werden kann.
Nach Satz~\ref{buch:integration:satz:homotopie} ergeben die
die beiden Wege denselben Wert des Integrals.
Der Wert des Integrals über den Weg $\gamma_1$ ist
\[
\oint_{\gamma_1}
\frac{dz}{z-z_0}
=
\int_0^1
\frac{2\pi i nre^{2\pi int}}{z_0+re^{2\pi i nt}-z_0}
\,dt
=
\int_0^1
2\pi i n
\,dt
=
2\pi i n.
\]
Aufgelöst nach $n$ folgt daraus
\[
n
=
\operatorname{ind}_\gamma(z_0)
=
\frac{1}{2\pi i}
\oint_{\gamma_1}\frac{dz}{z-z_0}
=
\frac{1}{2\pi i}
\oint_{\gamma}\frac{dz}{z-z_0}
\]
Damit ist der Satz bewiesen.
\end{proof}
